\section{Градиент}

\begin{frame}[c,fragile]{Частные производные и градиент}
  \begin{columns}[T]
    \begin{column}{0.52\textwidth}
      \begin{block}{Пример}
        \small
        \[f(x, y) = x^2 + xy + 2y^2\]
        Производная по $x$: $y$ считаем числом. По $y$: $x$ считаем числом.
        \[\frac{\partial f}{\partial x} = 2x + y, \qquad \frac{\partial f}{\partial y} = x + 4y\]
        \[\nabla f = \begin{pmatrix} 2x+y\\ x+4y\end{pmatrix},\qquad
          \nabla f(1, 1) = \begin{pmatrix} 3\\ 5\end{pmatrix}\]
      \end{block}
    \end{column}
    \begin{column}{0.44\textwidth}
      \centering
      \begin{tikzpicture}[scale=0.62]
        % level lines of x^2 + xy + 2y^2 = c
        \foreach \c in {1,2.5,4,6} {
          \draw[level line,rotate=-22.5] (0,0) ellipse ({sqrt(\c/0.793)} and {sqrt(\c/2.207)});
        }
        \fill[acDark] (1,1) circle (2.2pt);
        \draw[-{Stealth[length=2.2mm]},acOrange,line width=1.2pt] (1,1) -- ++(0.36,0.6)
          node[chart label,anchor=south,text=acOrange] {$\nabla f$};
        \draw[-{Stealth[length=2.2mm]},acDark,line width=1.2pt] (1,1) -- ++(-0.36,-0.6)
          node[chart label,anchor=north east,text=acDark] {$-\nabla f$};
        \fill[acDark] (0,0) circle (1.4pt);
      \end{tikzpicture}

      \smallskip
      \footnotesize\color{acGray}Линии уровня: $f$ одинакова вдоль каждой.
      Градиент перпендикулярен им и смотрит в сторону роста.
    \end{column}
  \end{columns}

  \medskip
  {\footnotesize\color{acGray}\href{https://www.geogebra.org/3d?command=f(x,y)=x^2%2Bx*y%2B2*y^2%3BSetFilling(f,0.35)%3BB=Point(f)%3BSetCoords(B,1,1,4)%3BA=(x(B),y(B),0)%3Bg=Vector(B,B%2B(2x(B)%2By(B),x(B)%2B4y(B),0))%3BSetColor(g,%22Orange%22)%3BSetLineThickness(g,10)}{\textcolor{acMid}{Откройте в GeoGebra 3D}}:
  точку $B$ можно двигать по поверхности, оранжевый вектор --- $\nabla f$ в этой точке.\par}
\end{frame}

\begin{frame}[c]{Градиент и сдвиг параметров}
  \begin{columns}[T]
    \begin{column}{0.47\textwidth}
      \begin{block}{Определение}
        \small
        Градиент --- вектор из всех частных производных:
        \[\nabla f(w) = \Bigl(\frac{\partial f}{\partial w_1}, \dots, \frac{\partial f}{\partial w_k}\Bigr)\]
      \end{block}
      \begin{block}{Сдвиг $h$}
        \small
        На сколько меняем каждый параметр: $w \to w + h$.
        Пример: $(1, 1) + (-0{,}15;\ -0{,}25) = (0{,}85;\ 0{,}75)$.
      \end{block}
    \end{column}
    \begin{column}{0.47\textwidth}
      \begin{block}{Линейное приближение}
        \small
        Если $h$ маленький, функция рядом почти плоская:
        \[f(w + h) \approx f(w) + \nabla f(w)\cdot h\]
        Справа скалярное произведение:
        \[\nabla f\cdot h = \frac{\partial f}{\partial w_1} h_1 + \dots + \frac{\partial f}{\partial w_k} h_k\]
        Каждый параметр сдвинули --- функция изменилась примерно на «наклон $\times$ сдвиг».
      \end{block}
    \end{column}
  \end{columns}
\end{frame}

\begin{frame}[c]{Почему шаг против градиента уменьшает функцию}
  \begin{columns}[T]
    \begin{column}{0.47\textwidth}
      \begin{block}{Куда шагнуть}
        \small
        Изменение функции:
        \[\nabla f\cdot h = \lVert\nabla f\rVert\,\lVert h\rVert\cos\varphi\]
        При шаге фиксированной длины оно самое отрицательное при $\cos\varphi = -1$,
        то есть когда $h$ смотрит \textbf{против} градиента.
      \end{block}
    \end{column}
    \begin{column}{0.47\textwidth}
      \begin{block}{Шаг ГД}
        \small
        Берём $h = -\eta\nabla f$:
        \[f(w - \eta\nabla f) \approx f(w) - \eta\lVert\nabla f\rVert^2 < f(w)\]
      \end{block}
      \begin{block}{Проверим на числах}
        \small
        $f = x^2 + xy + 2y^2$, $w = (1, 1)$, $f = 4$, $\nabla f = (3, 5)$, $\eta = 0{,}05$.
        Прогноз $4 - 0{,}05\cdot 34 = 2{,}3$, на деле $f(0{,}85;\ 0{,}75) \approx 2{,}49$.
      \end{block}
    \end{column}
  \end{columns}
\end{frame}

\begin{frame}[c]{Много параметров: линейная регрессия}
  \small
  Объект --- $k$ признаков, модель $\hat y_i = w_1 x_{i1} + \dots + w_k x_{ik}$. Ошибка
  $\;\mathrm{err}_i = \hat y_i - y_i$.

  \medskip
  \begin{columns}[T]
    \begin{column}{0.47\textwidth}
      \begin{block}{Шаг 1. Один объект}
        \small
        Цепное правило:
        \[\frac{\partial}{\partial w_j}\,\mathrm{err}_i^2 = 2\,\mathrm{err}_i\cdot x_{ij}\]
      \end{block}
    \end{column}
    \begin{column}{0.47\textwidth}
      \begin{block}{Шаг 2. Среднее по объектам}
        \small
        Производная суммы --- сумма производных:
        \[g_j = \frac{2}{n}\sum_{i=1}^{n} x_{ij}\,\mathrm{err}_i\]
      \end{block}
    \end{column}
  \end{columns}

  \medskip
  \lead{Это просто цикл:} для каждого параметра $j$ пробегаем объекты и копим сумму.
  Осталось заметить, что такой цикл уже есть в numpy.
\end{frame}

\begin{frame}[c]{Шаг 3. Цикл --- это $X^\top \mathrm{err}$}
  \begin{columns}[T]
    \begin{column}{0.50\textwidth}
      \small
      3 объекта, 2 признака:
      \[X = \begin{pmatrix} 1 & 2\\ 3 & 0\\ 0 & 1\end{pmatrix},\qquad
        \mathrm{err} = \begin{pmatrix} 1\\ -1\\ 2\end{pmatrix}\]
      \textbf{Циклом:}
      \begin{align*}
        \textstyle\sum_i x_{i1}\mathrm{err}_i &= 1\cdot1 + 3\cdot(-1) + 0\cdot2 = -2\\
        \textstyle\sum_i x_{i2}\mathrm{err}_i &= 2\cdot1 + 0\cdot(-1) + 1\cdot2 = 4
      \end{align*}
    \end{column}
    \begin{column}{0.46\textwidth}
      \small
      \textbf{Матрицей:} столбец $j$ матрицы $X$ --- это строка $j$ в $X^\top$.
      \[X^\top \mathrm{err} =
        \begin{pmatrix} 1 & 3 & 0\\ 2 & 0 & 1\end{pmatrix}
        \begin{pmatrix} 1\\ -1\\ 2\end{pmatrix} =
        \begin{pmatrix} -2\\ 4\end{pmatrix}\]
      \lead{Ответы совпали.}
      \[\nabla \mathrm{MSE} = \frac{2}{n}\,X^\top (Xw - y)\]
    \end{column}
  \end{columns}
\end{frame}

\begin{frame}[c]{Правило размерностей}
  \lead{Не помните, где транспонировать? Проверьте формы.}

  \medskip
  \centering
  \zebra{2}
  \begin{tabular}{lcl}
    \textbf{Что} & \textbf{Форма} & \textbf{Смысл}\\
    $X$ & $n\times k$ & объекты $\times$ признаки\\
    $w$ & $k\times 1$ & по весу на признак\\
    $Xw - y$ & $n\times 1$ & по ошибке на объект\\
    $\nabla$ & $k\times 1$ & \textbf{как у $w$}: по числу на параметр\\
  \end{tabular}

  \medskip
  \raggedright\small
  Из $X$ ($n\times k$) и ошибок ($n\times1$) получить $k\times1$ можно только так:
  $(k\times n)\cdot(n\times 1)$, то есть $X^\top\mathrm{err}$. Если формы не сходятся --- формула неверна.

  \medskip
  \centerline{\texttt{grad = 2 / n * X.T @ (X @ w - y)}}
\end{frame}
